Statement-only fresh-review packet. No printed proof or proof-scope metadata included.

\begin{proposition}[Every normal irreducible support is a smooth torus quotient]\label{cdp:st:C.2.UN.5}


Let $f:Y\to B$ be a proper holomorphic map from a smooth
complex threefold to a smooth curve, with connected fibres
and smooth general fibre a complex two-torus. Let $mS$ be
a full fibre whose reduced support is irreducible and normal.
Then $S$ is smooth and $e(S)=0$. More precisely, the
normalized base change $t=s^m$ gives a connected two-torus
$A$ and a finite unramified cyclic map
\[
                         A\longrightarrow S
\]
of degree $m$. Thus $S$ is a two-torus or a bielliptic
surface. In the bielliptic case the order $n$ of $K_S$ is
$2,3,4$, or $6$, and $n\mid m$; in particular a reduced
normal fibre is a torus.

Every two-torus occurs with every positive multiplicity.
Every bielliptic surface of canonical order $n$ occurs
with every multiplicity divisible by $n$. These
realizations can be made in smooth compact threefolds
over $\PP^1$.
\end{proposition}

Cited internal statement:
\begin{proposition}[Isolated singularities are excluded and suspension germs are transverse products]\label{cdp:st:A.5.17}


Let $f:Y\to C_0$ have the hypotheses of \ref{cdp:st:A.5.15}.
No local equation $f=u h^m$, with $u$ a holomorphic
unit and $m\ge1$, can have a reduced equation $h$
with an isolated critical point at the origin.
Suppose its local equation at a point of a fibre is
\[
             f=(x^a-g(y,z))^m,
             \qquad a\ge2,\quad m\ge1,             \tag{A.J.38}\label{cdp:eq:A.J.38}
\]
where $g\ne0$ and $g(0,0)=0$. Then the reduced
plane germ $\{g=0\}_{\rm red}$ is smooth. If
$b$ is its multiplicity in $\operatorname{div}(g)$,
there are local holomorphic coordinates in which
$f=(x^a-y^b)^m$.
Consequently, at every point where the full fibre
is a uniform multiple of its reduced support and
the reduced equation has a nonzero Hessian, its
critical surface germ is a constant
transverse $A_{b-1}$ product. A nonordinary change
of transverse type at a point with nonzero Hessian
is excluded.
\end{proposition}

Cited primary loci only (no target proof): Buchdahl 1999 Theorem11, printedp297, https://www.numdam.org/article/AIF_1999__49_1_287_0.pdf ; Catanese 2002 arXiv:math/0111245v2 Definition3.1/Theorems3.2--3.3 p9 and Theorem4.6 pp13--14, https://arxiv.org/pdf/math/0111245 .
